\section{Weekly Summary}

\begin{tcolorbox}[colback=yellow!10,colframe=black!50,title=AI-Generated Content]
\noindent The summary below was written by Claude (Anthropic) from that week's regression outputs and series descriptions. It has not been reviewed or fact-checked by a human, and should be read as an automated, informal recap -- not analysis.
\end{tcolorbox}

Welcome back to the weekly ritual where we let a computer grab two economic time series that have never met, force them into a linear regression, and then act shocked at the results. This week's contestants included California house prices, a discontinued corporate loan survey, carbon dioxide from distillate fuel, Japanese retail stock sentiment, and, in a genuinely inspired pairing, SNAP benefit recipients in Walker County, Alabama versus retail sales in Estonia. As always: these are random pairings, no causality is implied, and the correct emotional response to any "significant" finding here is a raised eyebrow, not a policy memo.

Monday opened with a bang that turned into a whimper. House prices in Oxnard-Thousand Oaks-Ventura versus discontinued financial-business investment posted a jaw-dropping R-squared of 0.87 in the raw regression, the kind of number that makes you want to call a press conference. Then we detrended it and the R-squared face-planted to 0.23. Translation: two things that both went up over several decades will look madly in love until you remind them that everything went up over several decades. Classic trend mirage, exactly the sort of thing this project exists to ridicule. The rest of the week was a graveyard of near-misses -- San Luis Obispo house prices versus a loan survey couldn't clear significance in either version (congratulations on the honesty), and the emerging-markets bond index versus piped gas prices had a gorgeous raw fit that collapsed into statistical mush once detrended, its p-value ballooning to a proud 0.32.

The actually interesting cases were the quiet ones that survived the detrending gauntlet. U.S. Assets versus Asian business confidence (Wednesday) stayed significant in both regressions, and the relationship even flipped sign after detrending, which is the kind of plot twist that suggests something real is wobbling around under the shared drift. Thursday's health-care "other separations" versus retail trade hires held up too, and honestly two labor-market flow measures having a genuine link is almost disappointingly plausible -- we come here for nonsense, not for things that make sense. And then Saturday's grand finale: SNAP recipients in one Alabama county versus Estonian retail sales stayed significant even after detrending, a correlation so gloriously unhinged that I refuse to theorize about it beyond noting that the universe has a sense of humor.

So the scoreboard: a couple of raw hits that dissolved the instant we asked them to prove themselves, and a handful of relationships that stubbornly refused to die. Do remember that "survived detrending" is a low bar for evidence and not a claim that Alabama food assistance is secretly steering Baltic shopping habits. It's a hobby, the correlations are spurious until proven otherwise, and proving otherwise is emphatically not the goal. See you next week for more randomly arranged marriages of convenience between innocent data series.
